\chapter{Estado da Arte}\label{cap:estado-arte}

Este capítulo revê o trabalho existente nas duas tarefas em que a dissertação se foca,
isto é, o reconhecimento do número de camisola e a atribuição da equipa de cada jogador,
ambas propostas no âmbito do desafio SoccerNet. Antes de as abordar em detalhe, importa
enquadrar cada tarefa no \emph{pipeline} geral da reconstrução do estado de jogo e
apresentar as métricas utilizadas, que servem tanto para avaliar o sistema completo como
cada tarefa isoladamente. Em vez de percorrer os estudos um a um, a revisão está
organizada por tema, para deixar claras as principais linhas de investigação e os seus
limites.

\section{Reconstrução do Estado de Jogo}\label{sec:sistemas_gsr}

O conceito de reconstrução do estado de jogo foi proposto por \citet{somers2024soccernet}
como uma tarefa de referência na área da análise computacional do desporto, mais
propriamente no futebol. Os autores definiram, em simultâneo, a concetualização do
problema, o \emph{dataset} SoccerNet-GSR e a arquitetura de referência. Em termos práticos, o
objetivo consiste em converter a informação visual de vários \emph{frames} de um vídeo monocular
numa projeção bidimensional do terreno de jogo (Figura~\ref{fig:gsr-minimap}). Esta
transformação exige não só a deteção e localização contínuas das linhas do campo, de forma
a enquadrar espacialmente cada interveniente, mas também a sua atribuição relativamente à
equipa, à função desempenhada e ao número da camisola.

\begin{figure}[htbp]
    \centering
    \includegraphics[width=0.8\textwidth]{gamestate-example.jpg}
    \caption{Exemplo de reconstrução do estado de jogo, com a posição, a equipa
    e o número de camisola projetados num mapa bidimensional do campo. Adaptado
    de \citet{somers2024soccernet}.}
    \label{fig:gsr-minimap}
\end{figure}

Como ponto de partida para o \emph{pipeline} apresentado, a \emph{baseline} do SoccerNet-GSR recorre ao \emph{framework} TrackLab. Esta ferramenta adota uma estrutura modular configurada para
coordenar, de forma sequencial, a deteção, a reidentificação, o \emph{tracking} dos jogadores e a
calibração da câmara (Figura~\ref{fig:gsr-pipeline}). Neste ecossistema, a etapa final de
identificação desdobra-se em três vertentes. O modelo PRTreID extrai características
associadas à identidade, à equipa e à função de cada jogador. A ferramenta MMOCR reconhece
os números nas camisolas. Por fim, um algoritmo de agrupamento (\emph{clustering}) opera sobre as
representações visuais para inferir a que equipa pertence cada atleta. A presente
dissertação foca-se criticamente nestes dois últimos componentes, uma vez que constituem
os principais pontos de estrangulamento (\emph{bottlenecks}) de desempenho na solução proposta
por \citet{somers2024soccernet}.


\begin{figure}[htbp]
    \centering
    \includegraphics[width=\textwidth]{gsr_baseline_architecture.jpg}
    \caption{Arquitetura de referência do SoccerNet-GSR. Adaptado
    de \citet{somers2024soccernet}.}
    \label{fig:gsr-pipeline}
\end{figure}

Esta mesma arquitetura modular acaba por revelar as suas próprias fragilidades. Como
demonstraram os relatórios das edições seguintes do desafio \citep{cioppa2024soccernet,
giancola2025soccernet}, qualquer ganho real na projeção do mapa tático (\emph{minimap})
depende diretamente da qualidade da identificação local, e não apenas da capacidade de
detetar ou seguir os atletas. O percurso da solução vencedora da edição de 2024
\citep{golovkin2025minimap} é, a este respeito, profundamente revelador. Ao elevar o
GS-HOTA de $23.36\%$ da solução de referência para $63.81\%$, a sua estratégia converge,
de forma independente, com as duas intuições centrais desta dissertação. A primeira é que o
reconhecimento do número ganha em afastar-se de abordagens de OCR genéricas e adotar
soluções que integram conhecimento específico do domínio. A segunda é que a atribuição da
equipa ganha uma nova robustez quando se combina a aparência visual com a consistência
temporal e o contexto espacial do jogador. Importa, no entanto, sublinhar uma distinção de
âmbito. Essa solução vencedora distribuiu os seus ganhos por vários módulos, da deteção à
calibração, ao \emph{tracking} e ao pós-processamento. A presente dissertação, pelo
contrário, concentra-se exclusivamente nas duas tarefas de identificação, mantendo a deteção,
o \emph{tracking} e a calibração tal como na \emph{baseline}, e a reidentificação com o
seu modelo intacto, apenas acrescida de um descritor de cor ao serviço da atribuição da
equipa. É precisamente neste espaço de
oportunidade, e com esta delimitação deliberada, que o trabalho se posiciona. Trata-se de
isolar e redesenhar os dois componentes mais frágeis, de modo a que qualquer melhoria
observada possa ser atribuída, sem ambiguidade, a essas alterações, e não a uma otimização
global do sistema.

Convém situar este trabalho no percurso, curto, da própria tarefa. A reconstrução do estado
de jogo foi disputada em duas edições do desafio. Na primeira, em 2024, participaram
dezasseis equipas com trinta e quatro submissões, e a solução vencedora elevou o GS-HOTA da
solução de referência de $23.36\%$ para $63.81\%$ \citep{cioppa2024soccernet}. Na segunda,
em 2025, catorze equipas apresentaram sessenta e seis submissões sobre uma solução de
referência atualizada, com $29.01\%$, e a vencedora obteve $63.90\%$, com a solução vencedora de 2024 em
segundo lugar, com o mesmo valor de $63.81\%$ \citep{giancola2025soccernet}. Dois traços deste percurso importam para o que se segue. O
primeiro é que o valor máximo praticamente não se moveu entre as duas edições, apesar da
mudança de paradigma na leitura do número, sinal de que a tarefa está longe de resolvida e de
que nenhum módulo isolado a desbloqueia. O segundo é que os
organizadores registam, na edição de 2025, três tendências transversais aos sistemas
melhor classificados, a saber, o recurso a características CLIP e OSNet para a
reidentificação, a função e a equipa, o alisamento temporal da calibração por fluxo ótico, e a
aplicação de modelos de visão e linguagem à leitura do número
\citep{giancola2025soccernet}. A edição de 2026 do desafio já não incluiu a tarefa, substituída
por cinco outras, entre as quais a deteção de ações centrada no jogador, que reutiliza a
equipa e o número de camisola como identificadores \citep{cioppa2026soccernet}. A comparação
com estas classificações é retomada no Capítulo~\ref{cap:resultados}, com as reservas que a
diferença entre o conjunto de teste público e o conjunto do desafio impõe.

\section{Reconhecimento de Número de Camisola}

O reconhecimento automático do número da camisola a partir de sequências de vídeo de câmara
única constitui um desafio complexo na área da visão por computador, estando estritamente
condicionado por restrições estruturais. Devido à distância operacional da câmara em
relação aos jogadores e ao relvado, a resolução e definição dos dígitos é frequentemente
reduzida. Adicionalmente, as dinâmicas de movimento dos jogadores resultam em oclusões
frequentes, deformações nos tecidos das camisolas e variações extremas de pose, fazendo
com que os números raramente se encontrem num plano frontal e legível. Para mitigar estas
dificuldades, as abordagens propostas na literatura dividem-se, genericamente, em duas
grandes vertentes metodológicas.

\subsection{Abordagens por OCR Local e Modelos de Pose}

Os primeiros trabalhos abordaram a tarefa como um problema de deteção e classificação
de dígitos. \citet{gerke2015soccer} foram dos primeiros a aplicar redes neuronais
convolucionais (\emph{Convolutional Neural Networks}, CNN) ao reconhecimento do número
de camisola, estabelecendo a viabilidade da abordagem. \citet{liu2019poseguided} deram
um passo importante com uma R-CNN guiada por pose (\emph{pose-guided} R-CNN). Substituíram
a \emph{Region Proposal Network} (RPN) tradicional por uma variante de três classes
(fundo, pessoa e dígito) e incorporaram os \emph{keypoints} do corpo humano no
refinamento das propostas, melhorando a localização dos dígitos mesmo em poses
dinâmicas. O modelo foi validado num conjunto de dados criado especificamente para o
efeito, com anotações detalhadas de caixas de pessoa e de dígito.

Esta linha foi consolidada por \citet{liu2022jede} com o JEDE, um detetor universal que
combina uma \emph{Feature Pyramid Network} (FPN) com uma \emph{Region Proposal Network}
(RPN) e um ramo guiado por pose (\emph{pose-guided branch}) para as propostas de dígitos.
Testado em futebol e basquetebol, superou métodos clássicos como o Faster R-CNN e o Mask
TextSpotter, sobretudo na precisão e revocação da deteção.

Mais recentemente, \citet{koshkina2024general} propuseram uma abordagem distinta, que
trata o reconhecimento do número como um caso de reconhecimento de texto em cena
(\emph{scene text recognition}). Em vez de detetar cada dígito isoladamente, encadeiam
um classificador de legibilidade (\emph{legibility classifier}), que descarta as
\emph{frames} onde o número não está legível, com um modelo de reconhecimento de texto,
agregando depois as leituras ao longo do \emph{tracklet}. Os autores disponibilizam ainda
um conjunto de dados público, com anotações de legibilidade, e avaliam o desempenho tanto
ao nível da imagem como do \emph{tracklet}. Esta estratégia, ao filtrar previamente as
imagens úteis e ao consolidar a decisão no tempo, aproxima-se da lógica adotada nesta
dissertação.

\subsection{Abordagens temporais e por \emph{keyframes}}

Uma linha de investigação distinta tira partido da dimensão temporal do vídeo, uma vez que
a mesma camisola é vista em muitos \emph{frames}, mas nem todos úteis.
\citet{balaji2023jersey} introduziram um módulo de identificação de \emph{keyframes} numa
rede espácio-temporal, concentrando assim a análise nos \emph{frames} que contêm mais
informação e classificando cada dígito separadamente através de uma função de perda
multitarefa (\emph{multi-task loss}). A estratégia lida bem com oclusões e movimento
rápido, mas tem uma limitação que os próprios autores reconhecem. Ao tratar cada número
como um todo, e não dígito a dígito, o modelo dificilmente reconhece combinações que não
tenham surgido durante o treino.

O trabalho seguinte da mesma equipa, \citet{balaji2024dmae}, refina a abordagem com
\emph{Domain-Guided Masked Autoencoders} (d-MAE), treinados com desfoque de movimento
para simular condições reais, e um descodificador \emph{transformer} temporal. Os ganhos
reportados são modestos mas ilustrativos, cerca de $1{.}39\%$ com o módulo de agrupamento
e $0{.}77\%$ na localização, e reforçam uma ideia central. Em baixa resolução, a cor por
si só é insuficiente, e importa extrair características profundas e robustas para detetar
dígitos pequenos.

\subsection{Métodos apresentados ao desafio SoccerNet}\label{sec:jnr-desafio}

As duas linhas anteriores descrevem a investigação publicada em conferências e revistas.
Existe, porém, um segundo corpo de trabalho que raramente chega a esses canais e que é
particularmente informativo para esta dissertação, o das soluções apresentadas às
sucessivas edições do desafio SoccerNet. Interessam aqui não pelas classificações que
obtiveram, mas pelas opções técnicas que revelam, e sobretudo pela mudança de paradigma
que documentam ao longo de três anos. Salvo indicação em contrário, as descrições que se
seguem resumem os textos que as próprias equipas submeteram aos relatórios oficiais
\citep{cioppa2023soccernet, cioppa2024soccernet, giancola2025soccernet}.

A tarefa foi disputada isoladamente na edição de 2023, sobre um conjunto de
\emph{tracklets} já recortados, com quinze equipas e 157 submissões. Importa esclarecer
que esses \emph{tracklets} provêm do conjunto SoccerNet-Tracking e que a tarefa estava
formulada como classificação, com uma classe adicional para o número não visível, pelo que
os valores dessa edição não são comparáveis com os aqui reportados. O que dela se retira é
o método. Segundo os organizadores, a generalidade das equipas adotou um mesmo
procedimento em três fases. Primeiro uma deteção de texto sobre o recorte do jogador, com
recurso a modelos como o DBNet++, o MMOCR, o DeepSolo ou o YOLO. Depois um reconhecimento
do texto localizado, por motores de OCR afinados para o efeito, tipicamente o PP-OCRv3 ou
o PaddleOCR. E por fim uma votação maioritária que consolida as leituras de todas as
imagens de um mesmo \emph{tracklet} numa decisão única. Duas observações se impõem. A
primeira é que a agregação temporal por votação, que também se adota neste trabalho, já
era prática corrente da comunidade em 2023, e não constitui portanto uma novidade. A
segunda é que toda esta família assenta em tratar o número como texto arbitrário, que é
exatamente a limitação que a Secção~\ref{sec:met-jnr} procura ultrapassar.

Na primeira edição da reconstrução do estado de jogo, em 2024, as equipas melhor
classificadas abandonaram o OCR genérico e passaram a treinar modelos próprios,
especializados no número de camisola. \citet{golovkin2025minimap} organizam este panorama
em duas categorias, a saber, as adaptações de técnicas de OCR de uso geral, que localizam
e depois leem o texto, e as soluções desenhadas para o domínio desportivo. Sobre as
primeiras notam que a exatidão se degrada com jogadores distantes, com oclusões e com
orientações laterais, e sobre as segundas que ganham exatidão à custa de exigirem treino
específico para cada conjunto de dados. A própria solução vencedora ilustra bem a segunda
categoria, pois prescinde de detetar cada dígito com uma caixa delimitadora e usa em seu
lugar uma rede residual modificada, que recebe recortes de $32\times32$ píxeis e possui
duas cabeças de classificação, uma que decide se existe um dígito inicial e qual, e outra
que prevê o segundo dígito. A estrutura permite reconhecer o número mesmo quando parte
dele está ocluída, e dispensa a anotação de caixas por dígito. A equipa classificada em
segundo lugar seguiu caminho semelhante, treinando um detetor dedicado de dígitos sobre
dados públicos e cerca de mil imagens anotadas de propósito, e uma terceira equipa
recorreu ao CLIP em duas fases, com um modelo a decidir se há número visível e outro a
identificá-lo, afinados sobre três mil caixas anotadas manualmente. O traço comum a todas
é o custo de anotação, uma vez que cada uma exigiu a construção de um conjunto de treino
próprio, o que representa um esforço incompatível com as condições em que esta dissertação
foi realizada.

A edição de 2025 marca uma mudança clara. Os organizadores registam a aplicação
bem-sucedida de modelos de visão e linguagem à leitura do número, e as duas soluções mais
relevantes nesse sentido recorreram ao LLaMA-3.2-Vision e ao Qwen2-VL-Instruct. A solução
vencedora, descrita em detalhe por \citet{oo2026broadcast2pitch}, formula a identificação
completa do atleta como uma tarefa de geração autorregressiva multimodal. Em vez de treinar
um classificador por atributo, entrega ao modelo o recorte do jogador acompanhado de uma
instrução em linguagem natural que lhe pede a função, o número e a cor do equipamento, e lê
a resposta como uma sequência de símbolos. Os autores justificam a escolha por três razões,
a saber, a generalização a atributos nunca vistos, o que torna o reconhecimento aberto, a
redução do esforço de anotação, uma vez que basta ajustar a instrução, e a possibilidade de
transpor a solução para outros desportos.

Este mesmo trabalho fornece a comparação mais direta de que há registo entre as três
estratégias, por as avaliar dentro de um único \emph{pipeline} e sobre o mesmo conjunto de
teste. A combinação de modelos especializados por atributo, com o PRTreID para a função, o
EasyOCR para o número e uma rede residual para a cor, obtém $18{.}11\%$ de GS-HOTA. A
substituição por um modelo de visão e linguagem eleva o valor para $60{.}13\%$ com o CLIP e
para $61{.}48\%$ com o LLaMA-3.2-Vision. Ainda que estes valores digam respeito ao sistema
completo, e não isoladamente à leitura do número, a distância entre a primeira linha e as
restantes é demasiado grande para ser atribuída a outra coisa. Em sentido contrário, uma
das equipas melhor classificadas manteve-se no reconhecimento de texto convencional,
isolando primeiro a região do tronco a partir dos pontos de interesse do ViTPose e
aplicando depois o PARSeq, o que mostra que a via clássica continua viável quando a região
entregue ao leitor é a correta.

Vários participantes convergem, por fim, num mesmo diagnóstico. A equipa da Universidade de
Tóquio encerra o seu relato reconhecendo que, tal como na solução de referência, a maior
dificuldade permaneceu no reconhecimento do número, e aponta os modelos multimodais como o
caminho a seguir. A demonstração mais completa é, contudo, a de \citet{oo2026broadcast2pitch},
que decompõem a métrica oficial por atributo dentro do seu próprio sistema. Avaliando apenas
a posição no campo obtêm $79{.}52\%$. Acrescentando um atributo de cada vez à posição, o
valor desce marginalmente para $79{.}20\%$ com a função e para $73{.}58\%$ com a equipa, mas
cai para $64{.}70\%$ com o número de camisola, e com os quatro atributos em conjunto fica em
$61{.}48\%$. Os autores concluem que a leitura do número é o principal estrangulamento do
sistema, e explicam-no pela natureza da transmissão televisiva,
onde os dígitos surgem desfocados, ocluídos ou visíveis em poucas \emph{frames} de cada
\emph{tracklet}. Trata-se de confirmação externa e quantificada da premissa em que assenta
esta dissertação, e que motiva a escolha dos dois módulos a intervencionar.

\section{Atribuição da equipa}

A atribuição da equipa consiste em decidir a que equipa pertence cada jogador, quase
sempre por análise de características visuais. As abordagens diferem sobretudo no tipo de
sinal em que assentam.

\subsection{Métodos baseados em cor}

A cor do equipamento é a característica de escolha mais direta para separar equipas.
\citet{naik2022deepplayer} propuseram o DeepPlayer-Track, que combina informação de cor
com a distância do cosseno num \emph{tracker} do tipo SORT, apoiado num detetor YOLOv4.
O método distingue jogadores de forma robusta mesmo sob oclusão, mas os próprios autores
identificam uma fragilidade reveladora. Quando a cor do equipamento do árbitro não é a
esperada, o sistema não o deteta e confunde-o com um jogador. É um sintoma claro do
limite das abordagens puramente cromáticas, pois, quando as cores se aproximam ou fogem
ao previsto, o sinal deixa de ser fiável.

\subsection{Métodos por aprendizagem e agrupamento}

Para ganhar robustez, outros trabalhos substituem a cor por características aprendidas por
redes neuronais. \citet{alhejaily2023team} segmentam primeiro o campo e atribuem as
equipas com \emph{autoencoders} convolucionais treinados de ponta a ponta sobre dados não
rotulados, sem conhecimento prévio das cores de cada equipa, o que lhes confere capacidade
de generalização entre jogos diferentes. Ainda sem depender da cor,
\citet{istasse2019associative} propuseram aprender \emph{embeddings} associativos que
aproximam, no espaço latente, os jogadores da mesma equipa e afastam os das equipas
opostas, dispensando rótulos de cor explícitos. É também a filosofia da solução de
referência do desafio, que agrupa os jogadores sobre as representações visuais geradas
pelo PRTreID \citep{somers2024soccernet}.

\subsection{Métodos apresentados ao desafio SoccerNet}\label{sec:equipa-desafio}

As duas subsecções anteriores abordam a atribuição da equipa como um problema de separar os
jogadores em dois grupos a partir da sua aparência. As soluções apresentadas ao desafio, por
atuarem dentro de um sistema completo e serem avaliadas pela métrica oficial, revelam que o
problema tem na verdade duas etapas distintas, e que a segunda tem sido tratada com muito
menos atenção do que a primeira.

\citet{golovkin2025minimap} sintetizam o panorama de forma útil. A generalidade dos sistemas
começa por extrair, para cada jogador, uma representação da sua aparência, seja um histograma
de cor, seja um vetor de reidentificação treinado por aprendizagem contrastiva. Essas
representações são depois agrupadas por métodos não supervisionados, tipicamente o
\emph{k-means} ou o DBSCAN, em dois a cinco grupos, consoante se pretenda separar apenas as
duas equipas ou também os árbitros e os guarda-redes. A alternativa a esta estrutura é a
aprendizagem de representações discriminativas da equipa, como a proposta por
\citet{istasse2019associative}, que dispensa o agrupamento por produzir diretamente vetores
que aproximam colegas de equipa e afastam adversários.

Na edição de 2024, a equipa vencedora substituiu o agrupamento genérico por um modelo
treinado para o efeito, uma OSNet ajustada a 111 equipamentos distintos, e passou a procurar
cinco grupos em vez de dois, a saber, as duas equipas, os árbitros e os dois guarda-redes. O
procedimento é revelador do cuidado que a tarefa exige, pois para estimar os grupos das
equipas consideram-se apenas os jogadores a menos de trinta metros do centro do campo, de
modo a excluir os guarda-redes, que são depois identificados separadamente dentro das áreas.
Na edição de 2025 predominam duas orientações, a delegação da tarefa no mesmo modelo
multimodal que lê o número de camisola, como faz a solução vencedora, que pede ao modelo a
cor do equipamento e agrupa depois por essa cor, e o reforço da informação de aparência,
como na solução que segmenta a silhueta do jogador e dela extrai um descritor com as cores do
equipamento e as respetivas áreas.

Merece registo, pelo que ensina, o percurso da equipa que ficou em último lugar em 2024.
Propôs resolver a função e a equipa globalmente para todo o vídeo, através de uma rede
siamesa, de um agrupamento por \emph{k-means} com vários valores de $k$ e de uma floresta
aleatória que escolhia o agrupamento mais plausível a partir de características de contexto
de jogo. O resultado ficou abaixo da solução de referência, e os autores atribuem-no ao facto
de terem treinado o modelo com posições da anotação e o terem avaliado sobre deteções reais.
É um exemplo claro de como a qualidade das etapas anteriores condiciona tudo o que vem
depois.

Separar os jogadores em dois grupos resolve, no entanto, apenas metade do problema. Como a
métrica oficial exige que cada jogador seja identificado como pertencendo à equipa da
esquerda ou à da direita, é ainda necessário decidir qual dos dois grupos ocupa cada lado, e
é aqui que as soluções mais divergem. A abordagem mais comum consiste em comparar a posição
horizontal média dos jogadores de cada grupo, atribuindo o lado direito àquele cuja média
está mais à direita, e é o que faz a solução vencedora de 2025, que assume explicitamente
que os jogadores tendem a permanecer do seu lado do campo. A solução vencedora de 2024
identifica, contudo, uma fragilidade desta suposição, pois substitui a média única por uma
votação, em que cada \emph{frame} vota no grupo que nela surge mais à direita, e justifica a
alteração afirmando que a média simples ao longo de todo o clipe pode falhar em sequências
curtas ou quando apenas parte do campo está visível. Uma terceira via recorre à posição do
guarda-redes, atribuindo essa função ao jogador mais afastado dentro de cada área e
submetendo depois o resultado a uma otimização que penaliza configurações impossíveis, como
a existência de dois guarda-redes do mesmo lado. Importa deixar claro, por honestidade de
atribuição, que o recurso ao guarda-redes como referência de lado não é original deste
trabalho, uma vez que já se encontra nesta solução e, de forma implícita, na vencedora de
2024, que deteta os guarda-redes separadamente.

Duas lacunas atravessam, ainda assim, todas estas soluções. A primeira é que nenhuma
quantifica com que frequência a suposição da posição média falha, pois a vencedora de 2024
afirma que falha sem indicar em quantos casos nem em que condições, e sem essa medida não é
possível saber se o problema é marginal ou estrutural. A segunda é que a etapa de atribuição
do lado nunca é avaliada isoladamente, diluindo-se o seu efeito na métrica global, que agrega
a localização, o \emph{tracking}, a função, o número e a equipa. O Capítulo~\ref{cap:resultados}
responde a ambas, medindo sobre os 160 vídeos das três divisões do conjunto de dados a
frequência com que o critério da posição média se inverte, e avaliando separadamente o
contributo de cada nível da cascata de decisão descrita na Secção~\ref{sec:mapeamento}.

\section{Métricas de Avaliação}\label{sec:metricas}

A avaliação da reconstrução do estado de jogo suscita desafios particulares. Localizar
corretamente cada jogador no campo é apenas parte do problema. Torna-se igualmente
fundamental identificar quem cada um efetivamente é, ou seja, determinar o número de
camisola, a equipa e a função no jogo. Para responder a esta dupla exigência, a métrica
oficial do desafio, o GS-HOTA, foi introduzida por \citet{somers2024soccernet} em
simultâneo com a própria tarefa. Derivada do HOTA \citep{luiten2021hota}, uma métrica
consolidada no \emph{tracking} de múltiplos objetos, o GS-HOTA estende o conceito original
ao incorporar a dimensão da identidade, articulando a precisão da localização espacial com
a exatidão dos atributos do jogador. A diferença em relação ao HOTA clássico é
determinante. Para que uma deteção seja considerada válida, o GS-HOTA exige a
correspondência integral de todos os atributos com a verdade de referência (\emph{ground
truth}). A falha num único atributo penaliza a deteção, o que torna esta métrica
particularmente rigorosa e sensível a erros de identificação.

Por ser uma métrica global, o GS-HOTA reflete o desempenho conjunto de todos os
componentes, da deteção ao \emph{tracking} e da calibração à identificação, o que
dificulta o isolamento do contributo de cada um. Uma forma de contornar esta limitação,
adotada por \citet{oo2026broadcast2pitch} e retomada no Capítulo~\ref{cap:resultados},
consiste em calcular a métrica com apenas parte dos atributos ativados, começando pela
posição e acrescentando a função, a equipa e o número, de modo que a diferença entre dois
cálculos consecutivos mede o custo de cada atributo. Para avaliar separadamente as duas
tarefas em estudo, recorre-se ainda a métricas complementares, largamente aceites na
literatura. A precisão, a revocação e o F1 medem a qualidade da classificação, isto é,
quantas leituras ou atribuições estão corretas e a que custo de omissões. A exatidão dá a
fração de acertos. A precisão média (mAP) aplica-se quando a tarefa é formulada como
deteção. A conjugação dos dois tipos de métricas, a oficial para o efeito no sistema
completo e as dedicadas para o desempenho isolado de cada sub-tarefa, permite uma
avaliação ao mesmo tempo comparável ao desafio e informativa quanto à contribuição
específica de cada melhoria.

\section{Síntese e Lacunas Identificadas}

A análise da literatura revela um padrão claro. A fragilidade do sinal visual de baixo
nível leva a generalidade dos trabalhos a incorporar mecanismos adicionais de extração de
características, tais como \emph{autoencoders}, arquiteturas de \emph{backbone} mais
robustas ou estratégias de agrupamento mais elaboradas. Observa-se igualmente uma
evolução incremental do estado da arte, na qual os métodos mais recentes procuram refinar
abordagens prévias, como atesta a transição das redes R-CNN guiadas por pose para o JEDE,
ou a sucessão do reconhecimento de \emph{keyframes} pelos \emph{Masked Autoencoders}.

Confrontando esta literatura com as soluções apresentadas ao desafio, identificam-se duas
lacunas que motivam o presente trabalho:

\begin{enumerate}
    \item \textbf{Desequilíbrio de foco.} A investigação na área tem priorizado as
    componentes de \emph{tracking} e deteção de jogadores, relegando a sua identificação,
    designadamente o número da camisola e a atribuição da equipa, para um plano
    secundário. Contudo, é precisamente nesta vertente que a própria solução de referência
    (\emph{baseline}) apresenta as suas maiores limitações \citep{somers2024soccernet},
    diagnóstico que várias equipas do desafio confirmaram de forma independente ao
    identificarem a leitura do número como a dificuldade que lhes restou.

    \item \textbf{Ausência de avaliação isolada dentro do \emph{pipeline} completo.} As
    duas famílias de trabalho falham em sentidos opostos. Os estudos académicos avaliam os
    seus métodos de forma isolada e em \emph{datasets} próprios, sem os submeterem a um
    \emph{pipeline} completo de reconstrução do estado de jogo nem à métrica oficial. As
    soluções do desafio, pelo contrário, são avaliadas no \emph{pipeline} completo e sob a
    métrica oficial, mas otimizam todos os módulos em simultâneo, o que impede saber quanto
    do ganho se deve à identificação e quanto se deve à deteção ou à calibração.
\end{enumerate}

É esta segunda lacuna que define a estratégia experimental adotada. Ao intervir unicamente
na leitura do número e na atribuição da equipa, e ao deixar a deteção, o \emph{tracking} e a
calibração exatamente como na \emph{baseline}, e a reidentificação com o seu modelo intacto,
qualquer variação de desempenho passa a ser atribuível sem ambiguidade ao que foi alterado. O trabalho
inscreve-se assim na mesma direção técnica das soluções mais competitivas, em particular no
recurso a modelos multimodais para a leitura do número, mas distingue-se pelo enfoque numa
implementação sob fortes restrições de recursos e por um desenho experimental que permite
isolar o contributo das duas tarefas em estudo.
